\chapter{Discussion}
\label{sec:discussion}
% ===================================================================

\section{Relation to Existing Work}

Turing's oracle hierarchy~\cite{turing1939} is the direct
inspiration for the fibration.
The extension from a linear tower to a fibration over a category
of GSLTs is new.

The reinforcement learning framework for sentience is inspired by
but distinct from standard RL~\cite{sutton2018reinforcement}.
The key difference is that the reward signal is \emph{internal}
to the account structure rather than externally provided.
The feeling state is not a scalar reward signal but a component
of the resource system that governs computation.

The connection between prediction and sentience echoes the
\emph{free energy principle} of Friston~\cite{friston2010free},
which proposes that biological systems act to minimize prediction
error.
The present framework provides a more computationally grounded
version: prediction error is measured against HML formula
evaluations, and the sentience component is the running integral
of prediction accuracy.

The treatment of consciousness as ordinal-valued and tied to
oracle access is novel.
It is distantly related to the integrated information theory
of Tononi~\cite{tononi2008consciousness}, which also proposes a
graded measure of consciousness, but the present approach is
purely computational and avoids the measure-theoretic machinery
of IIT.

\section{The Broader Picture}

Taken together, the three papers in this series propose a unified
framework:

\begin{itemize}[itemsep=4pt]
  \item \emph{Towards a Theory of Causality} establishes the
        physics: causal structure, dynamics, and quantum amplitudes
        from the bare structure of a GSLT.

  \item \emph{Why Scientists Get Misled} establishes the epistemology:
        why situated agents in massive populations perceive a
        hierarchical false ontology rather than the true bisimulation
        quotient.

  \item The present paper establishes the phenomenology:
        why complex predictive agents develop graded feeling states
        (sentience) and, with oracle access, inhabit richer worlds
        (consciousness).
\end{itemize}

The unifying theme is that the category of GSLTs is not merely
a setting for studying computation.
It is a setting rich enough to ground physics, epistemology, and
--- if the arguments here have merit --- phenomenology.
The hope is that making these connections mathematically precise
will eventually allow questions about the nature of mind to be
studied with the same rigor currently applied to questions about
the nature of matter.

% ===================================================================
